\section{Large-scale efforts in AI-assisted mathematics}\label{supp:practice}

This section gives the fuller account of OpenAI's work summarized in Section~\ref{sec:std-practice}, and records how
the accounts used in Section~\ref{sec:std-compare} were prepared.

The most visible examples are the releases of OpenAI in 2026, which illustrate several distinct ways of combining
AI discovery with verification. In February, OpenAI submitted proof attempts by an internal model to First Proof,
a set of ten research-level problems posed by experts in their fields \citep{openai_firstproof_2026}. The model
was run ``with limited human supervision'': the researchers sometimes suggested retrying strategies that had
appeared fruitful, asked the model to expand or clarify parts of a proof after expert feedback, facilitated a
back-and-forth between the model and ChatGPT ``for verification, formatting, and style'', and for some problems
presented the best of a few attempts selected by human judgment. One attempt first believed likely correct was
later judged incorrect in the light of the official commentary, OpenAI states that the process was not as clean as a
properly controlled evaluation, and the accompanying preprint includes an appendix of prompt patterns. In
August, OpenAI published ten results in mathematics and theoretical computer science, with certificates in the
proof assistant Lean \citep{openai_tenproofs_2026}. In September it announced a resolution of the Navier--Stokes
problem, with a written proof and a Lean formalization \citep{openai_navierstokes_2026}. OpenAI describes a
system of coordinating agents, powered by an internal model, with access to a cached copy of the internet and to
code execution, and subdivided into groups that communicate internally; the group that produced the result
involved of the order of ten thousand concurrent agents. Different groups were prompted with different variants of
the problem statement, two leading to a proof of regularity and two to a counterexample. A run on easier problems
produced, with nearly a hundred agents over about fifty hours, a result on the Euler equations, after which the
researchers moved agents to Navier--Stokes, prompted them with the Euler result, and switched them to a further
trained version of the model. Groups were encouraged to explore diverse approaches and were then
``cross-pollinated'' by using a coding agent to consolidate the most useful insights of each group into follow-up
prompts. The resolution came after about 88 hours and 2.7 million messages between agents, and the Lean
formalization took another 17 hours. In October OpenAI released about seven hundred
manuscripts produced by an internal model by one fixed procedure on some four thousand posed problems, together
with a Lean library covering a part of the results and a change log recording retractions and repairs
\citep{openai_math_2026}; this release did not use the large communicating configuration of the Navier--Stokes
run. The repository describes the manuscripts as ``raw research data, rather than polished academic papers'', and
states that unformalized results could contain errors. Later in October, OpenAI announced three further results, in
harmonic analysis, complexity theory and quantum field theory, again described its manuscripts as raw research
data, and tested three ways of turning them into understanding: releasing the materials and leaving researchers to
organize themselves; forming an external group of collaborators while releasing the materials; and forming an
external group and leaving the release to it \citep{openai_academic_2026}. The collaborations follow four stated
principles: transparency, no academic credit claimed by OpenAI for discoveries of its models, recognition of the
groups for the understanding they develop, and independence of the groups. OpenAI notes that the choice of
coordinators gives some researchers earlier access to the materials than others.

\begin{discussion}[the design of the panel]
When the author decided to make the use of AI central to the paper, the panel was described to the main assistant
as organized in the manner of the recent multi-agent efforts in mathematics, and the paper was conceived as having
two parts: the mathematical question, and the demonstration, open to inspection, of how AI was used to answer it.
To check this description, two accounts of OpenAI's methods were prepared independently: one by an agent of the
main assistant, which read OpenAI's public repositories, the arXiv audits, the AGMAI guidelines and press reports
but could not retrieve OpenAI's own pages, and one by ChatGPT, at the author's request, which cited OpenAI's posts
directly. The author then retrieved OpenAI's Navier--Stokes and First Proof posts and its October post on academic
collaborations in a browser, and the main assistant
compared them with both accounts; the posts confirm the facts that both accounts reported, and the text above
quotes them directly. The two accounts agree on
the facts used here. Both find that the resemblance holds for the principles
listed in Section~\ref{sec:std-compare}, that OpenAI used model-to-model checking and coordinated groups of agents, and that no published
record specifies dedicated reviewer or judge roles; the comparison in Table~\ref{tab:compare} is stated
accordingly. The use of a second provider to cross-check a description of the literature is an instance of the
guardrail of Section~\ref{sec:collab-guard}.
\end{discussion}
